\section{Лекция 10. Особые решения дифференциальных уравнений и дискриминантная кривая}
\noeqref{eq:l10_cauchy_res, eq:l10_cauchy_imp, eq:l10_loc_uniq}
\noeqref{eq:l10_disc_sys, eq:l10_ex_ode, eq:l10_ex_disc_sys}
\noeqref{eq:l10_ex_param_y, eq:l10_ex_diff, eq:l10_ex_gen_param}
\noeqref{eq:l10_ex_gen_exp, eq:l10_ex_tangency}
\noeqref{eq:l10_alg_disc, eq:l10_alg_osgood, eq:l10_alg_tangency}

\begin{flushright}
    27.10.26
\end{flushright}

\subsection{Понятие особого решения для разрешённых и не разрешённых относительно производной уравнений}

Сравним постановку задачи о поиске особых решений для двух основных типов
обыкновенных дифференциальных уравнений 1-го порядка, изученных в предыдущих
лекциях.

\subsubsection{Случай 1. Уравнение, разрешённое относительно производной}

Рассмотрим классическую задачу Коши (см. Лекции 1–3):
\begin{equation} \label{eq:l10_cauchy_res}
    \begin{cases}
        y' = f(x, y), \\
        y(x_0) = y_0.
    \end{cases} \tag{1}
\end{equation}

Напомним определение 2.2 из Лекции 2: решение $y = \phi(x)$ уравнения
$y' = f(x, y)$ называется \textdef{особым решением}, если через каждую его
точку проходит другое решение, не совпадающее с $\phi(x)$ ни в какой сколь
угодно малой окрестности этой точки.

\begin{remark}
    В случае уравнения $y' = f(x, y)$ наклон касательной в точке $(x_0, y_0)$
    однозначно задаётся самой правой частью $y'(x_0) = f(x_0, y_0)$, поэтому
    любые два решения, проходящие через общую точку $(x_0, y_0)$, автоматически
    имеют в ней общую касательную. Как было показано в Лекциях 2–4,
    единственность решения задачи Коши \eqref{eq:l10_cauchy_res} может
    нарушаться только в тех точках, где функция $f(x, y)$ \textbf{не
    удовлетворяет условию Липшица} по переменной $y$ (однако само по себе
    нарушение условия Липшица ещё не гарантирует потерю единственности, что
    демонстрирует критерий Осгуда для автономных уравнений из Лекции 4).
\end{remark}

\subsubsection{Случай 2. Уравнение, не разрешённое относительно производной}

Рассмотрим теперь задачу Коши для уравнения, не разрешённого относительно
производной (см. Лекцию 9):
\begin{equation} \label{eq:l10_cauchy_imp}
    \begin{cases}
        f(x, y, y') = 0, \\
        y(x_0) = y_0, \\
        y'(x_0) = y_1,
    \end{cases} \tag{2}
\end{equation}
где начальная тройка $(x_0, y_0, y_1)$ удовлетворяет уравнению
$f(x_0, y_0, y_1) = 0$.

Поскольку через одну точку $(x_0, y_0)$ могут проходить интегральные кривые с
разными угловыми коэффициентами касательных $y'(x_0)$ (что само по себе не
нарушает единственности решения задачи Коши \eqref{eq:l10_cauchy_imp}), в
определение особого решения необходимо включить требование \textbf{касания}.

\begin{definition}
    Решение $y = \phi(x)$ дифференциального уравнения $f(x, y, y') = 0$
    называется \textdef{особым решением}, если через каждую точку его графика
    проходит хотя бы ещё одно решение этого уравнения, \textdef{касающееся}
    кривой $y = \phi(x)$ в данной точке (то есть имеющее в ней ту же
    производную) и не совпадающее с $\phi(x)$ ни в какой окрестности этой
    точки.
\end{definition}

Геометрически особое решение $y = \phi(x)$ является \textdef{огибающей}
однопараметрического семейства интегральных кривых общего решения: в каждой
своей точке огибающая касается некоторой кривой из общего семейства.

\begin{figure}[H]
    \centering
    \begin{tikzpicture}[scale=1.2]
        % Оси координат
        \draw[->] (-2.8, 0) -- (2.8, 0) node[right] {$x$};
        \draw[->] (0, -0.5) -- (0, 2.8) node[above] {$y$};

        % Огибающая (особое решение y = 0.25 x^2)
        \draw[red, very thick, domain=-2.5:2.5, samples=80]
            plot (\x, {0.25*\x*\x})
            node[above left, red] {$y = \phi(x)$ (огибающая)};

        % Семейство кривых, касающихся огибающей в разных точках
        \foreach \xO in {-1.6, -0.8, 0, 0.8, 1.6} {
            \pgfmathsetmacro{\yO}{0.25*\xO*\xO}
            \pgfmathsetmacro{\slp}{0.5*\xO}
            \draw[blue, thick, domain=\xO-0.75:\xO+0.75, samples=40]
                plot (\x, {\yO + \slp*(\x - \xO) + 0.9*(\x - \xO)*(\x - \xO)});
            \fill[black] (\xO, \yO) circle (1.5pt);
        }
        \node[blue, right] at (1.3, 0.5) {частные решения семейства};
    \end{tikzpicture}
    \caption{Особое решение $y = \phi(x)$ уравнения $f(x, y, y') = 0$ как
    огибающая семейства интегральных кривых: в каждой точке касания оба
    решения имеют общую координату $y(x_0) = y_0$ и общую производную
    $y'(x_0) = y_1$.}
\end{figure}

\subsection{Дискриминантная кривая и алгоритм поиска особых решений}

В Лекции 9 (Теорема 9.2) было установлено, что если $f \in C^1(W)$ и в
начальной точке $(x_0, y_0, y_1)$ выполняются условия:
\begin{equation} \label{eq:l10_loc_uniq}
    \begin{cases}
        f(x_0, y_0, y_1) = 0, \\
        f'_{y'}(x_0, y_0, y_1) \neq 0,
    \end{cases} \tag{3}
\end{equation}
то задача Коши \eqref{eq:l10_cauchy_imp} имеет в некоторой окрестности точки
$x_0$ единственное решение.

Следовательно, вдоль любого особого решения $y = \phi(x)$ (где единственность
задачи Коши \eqref{eq:l10_cauchy_imp} нарушается в каждой точке) второе условие
в \eqref{eq:l10_loc_uniq} обязано нарушаться, то есть частная производная
$f'_{y'}$ должна обращаться в нуль!

\begin{theorem}[\textbf{Необходимое условие особого решения}]
    Пусть $f(x, y, y') \in C^1(W)$. Если гладкая функция $y = \phi(x)$ является
    особым решением уравнения $f(x, y, y') = 0$, то в каждой точке её графика
    при $y' = \phi'(x)$ одновременно выполняются два равенства:
    \begin{equation} \label{eq:l10_disc_sys}
        \begin{cases}
            f(x, y, y') = 0, \\[1mm]
            f'_{y'}(x, y, y') = 0.
        \end{cases} \tag{4}
    \end{equation}
\end{theorem}

\begin{definition}
    Кривая на плоскости $Oxy$, определяемая уравнением $\Phi(x, y) = 0$ или же
    $y = \phi(x)$ (то есть зависит от $x$, $y$ и не зависит от $y'$. Получают,
    выражая $y'$ из 2-го уравнения $f'_{y'}(x, y, y') = 0$ и подставляя
    в 1-ое уравнение $f(x, y, y') = 0$),
    называется \textdef{дискриминантной кривой} дифференциального
    уравнения $f(x, y, y') = 0$.
\end{definition}

\begin{remark}[\textbf{Почему дискриминантная кривая — лишь «подозрение» на особое решение}]
    Система \eqref{eq:l10_disc_sys} даёт только \textbf{необходимое}, но не
    достаточное условие особого решения. Поэтому всякая найденная
    дискриминантная кривая $y = \phi(x)$ требует обязательной двухэтапной
    проверки:
    \begin{enumerate}
        \item \textbf{Проверка того, что кривая является решением ОДУ.}
        Необходимо подставить функцию $y = \phi(x)$ и её производную
        $y' = \phi'(x)$ в исходное уравнение $f(x, y, y') = 0$ (иногда вдоль
        дискриминантной кривой наклон касательной $\phi'(x)$ не совпадает со
        значением $y'$, при котором зануляется $f'_{y'}$, и тогда кривая
        вообще не является решением).
        \item \textbf{Проверка нарушения единственности (касания с другими
        решениями).} Необходимо убедиться, что через каждую точку кривой
        $y = \phi(x)$ проходит другое решение из общего семейства, касающееся
        $\phi(x)$ в этой точке и не совпадающее с $\phi(x)$ в её окрестности.
    \end{enumerate}
\end{remark}

\subsection{Пример нахождения особого решения}

\begin{example}
    Найти общее и особое решения дифференциального уравнения:
    \begin{equation} \label{eq:l10_ex_ode}
        -y + x + y' - \ln y' = 0 \quad (y' > 0). \tag{5}
    \end{equation}

    \begin{proof}[Решение]~\\
        \begin{quote}
            \textbf{Этап 1. Поиск дискриминантной кривой (подозрение на особое
            решение).}\\
            Обозначим левую часть уравнения через
            \[
                F(x, y, y') = x - y + y' - \ln y'.
            \]
            Вычислим частную производную по $y'$:
            \[
                F'_{y'}(x, y, y') = 1 - \frac{1}{y'}.
            \]
            Составим систему \eqref{eq:l10_disc_sys} для нахождения
            дискриминантной кривой:
            \begin{equation} \label{eq:l10_ex_disc_sys}
                \begin{cases}
                    x - y + y' - \ln y' = 0, \\[1mm]
                    1 - \dfrac{1}{y'} = 0.
                \end{cases} \tag{6}
            \end{equation}
            Из второго уравнения системы находим $y' = 1$. Подставляя $y' = 1$
            в первое уравнение системы \eqref{eq:l10_ex_disc_sys}, исключаем
            $y'$ и получаем уравнение дискриминантной кривой:
            \[
                x - y + 1 - \ln 1 = 0 \implies y = x + 1.
            \]
            Проверим, является ли прямая $y = x + 1$ решением исходного
            уравнения \eqref{eq:l10_ex_ode}: её производная равна
            $y' = (x + 1)' = 1$, и при подстановке в \eqref{eq:l10_ex_ode}
            получаем верное тождество:
            \[
                -(x + 1) + x + 1 - \ln 1 = 0 \iff 0 = 0.
            \]
            Следовательно, прямая $y = x + 1$ — решение уравнения (подозрение
            на особое решение).

            \textbf{Этап 2. Нахождение общего решения методом введения
            параметра.}\\
            Разрешим уравнение \eqref{eq:l10_ex_ode} относительно $y$ и введём
            параметр $p = y' > 0$ ($dy = p \, dx$):
            \begin{equation} \label{eq:l10_ex_param_y}
                y = x + p - \ln p. \tag{7}
            \end{equation}
            Возьмём полный дифференциал от обеих частей равенства
            \eqref{eq:l10_ex_param_y} и заменим $dy$ на $p \, dx$:
            \[
                dy = dx + dp - \frac{dp}{p} \implies
                p \, dx = dx + dp - \frac{dp}{p}.
            \]
            Аккуратно сгруппируем слагаемые с $dx$ и слагаемые с $dp$:
            \begin{equation} \label{eq:l10_ex_diff}
                (1 - p) \, dx + \left(1 - \frac{1}{p}\right) dp = 0 \iff
                (p - 1)\left(\frac{dp}{p} - dx\right) = 0. \tag{8}
            \end{equation}
            Равенство \eqref{eq:l10_ex_diff} распадается на два случая:
            \begin{enumerate}
                \item $p - 1 = 0 \implies p = 1$. Подстановка в
                \eqref{eq:l10_ex_param_y} вновь даёт уже найденную нами
                дискриминантную прямую $y = x + 1$.
                \item $\frac{dp}{p} - dx = 0 \iff dx = \frac{dp}{p}$.
                Интегрируя это уравнение (с учётом условия $p > 0$), находим
                выражение для $x$:
                \[
                    x = \ln p + C, \quad C \in \R.
                \]
                Подставим найденное $x = \ln p + C$ в уравнение
                \eqref{eq:l10_ex_param_y}:
                \[
                    y = (\ln p + C) + p - \ln p = p + C.
                \]
            \end{enumerate}
            Таким образом, общее решение уравнения в параметрической форме
            имеет вид:
            \begin{equation} \label{eq:l10_ex_gen_param}
                \begin{cases}
                    x = \ln p + C, \\
                    y = p + C,
                \end{cases}
                \quad p > 0, \; C \in \R. \tag{9}
            \end{equation}
            Исключая параметр $p = e^{x - C}$ из системы
            \eqref{eq:l10_ex_gen_param}, получаем явное представление общего
            решения в виде семейства экспонент:
            \begin{equation} \label{eq:l10_ex_gen_exp}
                y(x) = e^{x - C} + C. \tag{10}
            \end{equation}

            \textbf{Этап 3. Строгая проверка того, что $y = x + 1$ — особое
            решение.}\\
            Выясним, существуют ли такие значения параметра $p > 0$ и константы
            $C \in \R$, при которых кривая из общего семейства
            \eqref{eq:l10_ex_gen_param} проходит через заданную точку
            $(x_0, x_0 + 1)$ прямой $y = x + 1$ и касается её в этой точке:
            \begin{equation} \label{eq:l10_ex_tangency}
                \begin{cases}
                    y = x + 1, \\
                    y' = 1.
                \end{cases} \tag{11}
            \end{equation}
            Подставляя координаты \eqref{eq:l10_ex_gen_param} в первое
            условие $y = x + 1$, получаем:
            \[
                p + C = (\ln p + C) + 1 \implies p = \ln p + 1 \implies p = 1,
            \]
            что полностью согласуется с условием $y' = 1$! При $p = 1$ точка
            касания кривой семейства \eqref{eq:l10_ex_gen_param} с прямой
            $y = x + 1$ имеет координаты:
            \[
                \begin{cases}
                    x = \ln 1 + C = C, \\
                    y = 1 + C.
                \end{cases}
            \]
            Значит, через \textbf{любую} наперёд заданную точку $(x_0, x_0 + 1)$
            прямой $y = x + 1$ проходит экспонента $y = e^{x - x_0} + x_0$ из
            общего семейства (при выборе $C = x_0$), которая касается прямой
            $y = x + 1$ в точке $x = x_0$ и не совпадает с ней ни в какой
            окрестности точки $x_0$ (так как $e^{x - x_0} + x_0 > x + 1$ при
            всех $x \neq x_0$). Следовательно, прямая $y = x + 1$ является
            \textdef{особым решением}.

            \textbf{Ответ:} общее решение $y = e^{x - C} + C$, особое решение
            $y = x + 1$.
        \end{quote}
    \end{proof}
\end{example}

\begin{figure}[H]
    \centering
    \begin{tikzpicture}[scale=1.1]
        % Оси координат
        \draw[->] (-2.8, 0) -- (3.0, 0) node[right] {$x$};
        \draw[->] (0, -1.6) -- (0, 3.6) node[above] {$y$};

        % Особое решение (огибающая) y = x + 1
        \draw[red, very thick] (-2.4, -1.4) -- (2.4, 3.4)
            node[above right, red] {$y = x + 1$ (особое решение)};

        % Семейство экспонент y = e^{x - C} + C
        \foreach \c/\lbl in {-1.2/C_1, 0/C_2, 1.2/C_3} {
            \draw[blue, thick, domain=-2.5:\c+1.1, samples=70]
                plot (\x, {exp(\x - \c) + \c});
            \fill[black] (\c, \c + 1) circle (1.8pt);
            \draw[dashed, gray] (\c, 0) node[below, black] {\footnotesize $\lbl$}
                -- (\c, \c + 1);
        }
        \node[blue, above left] at (1.2, 2.9) {$y = e^{x - C} + C$};
    \end{tikzpicture}
    \caption{Семейство интегральных кривых $y = e^{x - C} + C$ уравнения
    $-y + x + y' - \ln y' = 0$ и их огибающая — особое решение $y = x + 1$,
    касающееся каждой экспоненты в точке $(C_i, C_i + 1)$.}
\end{figure}